% RESEARCH_DRAFT_SHA256: 287b79ce17feff211b604aec6ab35a88766971f57f6033290106836e59bd5cc2
% Layout and source-reference wrappers expanded for the web renderer.
\documentclass[11pt]{article}
\usepackage[margin=25mm]{geometry}
\usepackage{amsthm}
\usepackage{fontspec,unicode-math}
\setmainfont{Latin Modern Roman}
\setmathfont{Latin Modern Math}
\usepackage{amsmath,amssymb,mathtools}
\usepackage{xurl,hyperref}
\hypersetup{colorlinks=true,urlcolor=blue,linkcolor=blue}
\setcounter{secnumdepth}{0}
\setlength{\parindent}{0pt}
\setlength{\parskip}{5pt}
\setlength{\emergencystretch}{3em}
\newcommand{\sref}[2]{\hyperlink{source:#1:#2}{[S#2]}}
\newcommand{\cataloguescope}[1]{\paragraph{Source question.}#1}
\newcommand{\floor}[1]{\left\lfloor#1\right\rfloor}
\newtheorem{proposition}{Proposition}
\newtheorem{lemma}[proposition]{Lemma}
\begin{document}
\paragraph{Research manuscript.} OpenAI \texttt{gpt-6-astra}, reasoning \texttt{ultra}; 11 October 2026.
\paragraph{Preserved source body.} The complete catalogue mathematical text and sources follow with layout and source links expanded, including source qualifications and its historical preparation notes. New research begins at the section headed Research attempt.

% UnsolvedMath ID 30000602; source catalogue code OWR-1386-003
\section{Sampling Algorithms for Maximum-Cut Density}\label{30000602}
{\small UnsolvedMath ID 30000602 · Computer Science}
\subsection{Definitions and mathematical statement}

\paragraph{Shared notation.}Standard number systems are denoted by $\mathbb N,\mathbb Z,\mathbb Q,\mathbb R,\mathbb C$; index ranges, computational models and specialized notation follow the cited problem source and the definitions below.


The maximum-cut size is $\max_{S\subseteq V}|E(S,V\setminus S)|$. A uniformly random induced subgraph is obtained by sampling a vertex subset of the specified size and retaining every edge between sampled vertices.
\paragraph{Statement recorded in the supplied source.}
For a finite graph $G$, determine the density of its maximum cut, or estimate it to prescribed accuracy from a uniformly random induced subgraph of fixed size. \href{https://doi.org/10.4171/owr/2006/48}{[S2]}
\paragraph{Scope qualification.}
The normalization of density, required accuracy, confidence level, and dependence of sample size on accuracy are not specified.
\paragraph{Source status.}
The supplied catalogue labels this record open. Its source-status metadata was checked against the supplied snapshot on 2026-10-11. This import does not independently establish current open status or mathematical correctness. The catalogue literature assessment is dated 2026-08-21.
\subsection{Short English statement}
Estimate maximum-cut density from a small uniformly sampled induced subgraph.
\subsection{Sources}
\begin{itemize}
\item[S1] Ulam AI and the UnsolvedMath Contributors, UnsolvedMath: A Curated Collection of Open Mathematics Problems, record 30000602 (OWR-1386-003). Catalogue curation and metadata: CC BY 4.0. The supplied source question is retained; definitions and the English summary are editorial additions.
\url{https://huggingface.co/datasets/ulamai/UnsolvedMath}.
\item[S2] Combinatorics, Probability and Computing (2007), Oberwolfach Reports, DOI 10.4171/owr/2006/48, pp. 2875--2952
\url{https://doi.org/10.4171/owr/2006/48}.
\item[S3] graph limits and parameter testing
\url{https://users.renyi.hu/~sos/2006_Graph_limits_and_parameter_testing.pdf}.
\textit{Underlying Oberwolfach Reports material retains its source-specific licence, including CC BY-SA 4.0 where applicable; see THIRD\_PARTY\_NOTICES.md.}
\end{itemize}
\paragraph{Preparation and review.}
Original question attributed to the cited source and UnsolvedMath contributors. Codex (GPT-6; exact serving revision not exposed) assisted with definitions, English summaries, source inspection and formatting. No new human mathematical authorship or independent proof verification is asserted. Canonical ID 30000602 and source code OWR-1386-003 checked against the complete supplied registry and existing repository statements; known flagged repetitions and obvious identical conjectures excluded. Full mathematical-equivalence/literature review has not been performed. Statement preparation and final review followed contribution prompts C and D; automated checks verify structure and identity, not open status.
\label{end:30000602}

\section{Research attempt}
\subsection{Literature, source context and scope}
The full original catalogue body and source qualifications are preserved above. The relevant primary workshop passages were read in the cached publisher report. Historical claims and subsequent attributed results are distinguished from the mathematical work below. A restricted example, standard lemma or known theorem is not counted as new progress on the entire original target. Missing assumptions and repeated or rhetorical extracts are not silently turned into different solved problems.

\subsection{Why unbiased cut sampling does not survive optimization}
The report already discusses graph-parameter testability and a cut-distance sampling theorem. Its maximum-cut example is not treated here as an untouched conjecture. We examine one tempting elementary proof route. For a graph $G$ on $n\ge q\ge2$ vertices and a fixed cut $A\subseteq V$, define its density as $c_G(A)/\binom n2$, where $c_G(A)$ counts crossing edges. Sample a uniform $q$-vertex set $S$ and use the same normalization on $G[S]$.
\begin{proposition}
For every fixed cut $A$, the sampled density $c_{G[S]}(A\cap S)/\binom q2$ is unbiased for $c_G(A)/\binom n2$. But if $\mu(G)$ is the maximum cut density, then $\mathbb E\mu(G[S])\ge\mu(G)$, and equality need not hold.
\end{proposition}
\begin{proof}
Each unordered vertex pair is included in $S$ with probability $\binom q2/\binom n2$, by counting sets containing that pair. Summing the corresponding indicators over the fixed crossing edges and dividing by $\binom q2$ proves unbiasedness. Choose a maximizing cut $A_*$ in $G$. For every $S$, the maximum in $G[S]$ is at least the value of the restricted cut $A_*\cap S$. Taking expectations proves the inequality. For $G=K_3$ and $q=2$, every sample is an edge with maximum density 1, while the full graph's maximum cut contains two of its three pairs and has density $2/3$. The inequality is strict.
\end{proof}
\subsection{Verification and remaining gaps}
All three two-vertex samples of $K_3$ were enumerated, including the restrictions of a fixed optimal full-graph cut. Their fixed-cut mean is $2/3$, while their optimized mean is 1. Thus exchanging maximization and expectation is exactly the unjustified step in the proposed unbiased-estimator proof. This does not refute constant-size approximation with an appropriately larger sample, and it does not contradict the primary report's testing results. No dimension-independent sample-size/error guarantee is derived in this attempt.

\section{References and source attribution}
Primary workshop report: \url{https://doi.org/10.4171/owr/2006/48}. Relevant passages read: 2006/48 pp. 2891--2892, testable graph parameters and maximum-cut sampling.

All retrievals were checked on 11 October 2026. Original question attribution remains in the preserved Sources section. External theorems are attributed and not claimed as new discoveries or fully verified merely by citation.


\paragraph{Integration review.} Codex (GPT-6; exact serving revision not exposed) reviewed the written arguments and expanded layout and source-reference wrappers. Repository metadata and the production web renderer were checked using Python 3.13 and JavaScript with MathJax 3.2.2. This model-assisted review is not independent human proof verification; no PDF compilation or formal proof checking is asserted.

\section{Completion Estimate}
\noindent\textbf{COMPLETION: 0\%}
No substantiated new progress toward the entire original target is claimed. The worked results and failed route are delimited above; a known result, illustrative model, or auxiliary obstruction is not counted as original-target completion.

\section{Final status and attribution}
\textbf{Attempt status: unresolved.} This describes this research attempt, not a new ruling on the source catalogue's status or on external literature claims.
The mathematics and this self-audit were generated by OpenAI \texttt{gpt-6-astra}, reasoning \texttt{ultra}, on 11 October 2026. Codex (GPT-6; exact serving revision not exposed) orchestrated the batch and contributed formatting/repository review. No human mathematical authorship, independent human proof checking, or formal verification is claimed. The historical preparation paragraphs in the reproduced catalogue body concern that body only.

Contribution prompt A was applied to the complete definition and supplied source record, followed by the adversarial checks of prompt B. Local reproduction files are \texttt{batch016/checks.py}, \texttt{computations.json}, and the per-ID audit JSON. They were not public links at preparation time. Numerical checks are finite and do not prove unrestricted results. No LaTeX engine was found locally; compilation was not performed by this research agent.

\end{document}
